\section{Research Objectives}
\label{sec:research-objectives}

The primary objective of this research is to develop a precise localisation system for autonomous vineyard robots that enables accurate spatial alignment between the robot, the 3D models, and the environment.

The research pursues the following sub-objectives:

\begin{enumerate}

    \item Develop a LiDAR-inertial SLAM system that achieves sub-decimetre absolute trajectory error in GNSS-degraded vineyard environments, addressing the structural ambiguity of row-crop corridors.

    \item Design a precise-alignment pipeline that refines the SLAM pose to sub-centimetre accuracy by registering live sensor data against a 3D Gaussian Splatting reference model using point cloud registration.

    \item Evaluate two sensing modalities (stereo vision and LiDAR) for the precise-alignment stage, comparing their registration accuracy and failure modes under vineyard field conditions.

    \item Integrate both localisation stages into a hierarchical architecture that provides pose estimates suitable for downstream robotic manipulation, such that the corrected pose feeds directly into the pruning arm's path planner.

\end{enumerate}

\subsection{Contributions}
\label{subsec:contributions}

This thesis makes the following contributions:

\begin{enumerate}
    \item A LiDAR-inertial SLAM system augmented with sparse GNSS factor injection (GLIM-GNSS) that achieves 71--174\,mm absolute trajectory error across five vineyard datasets while utilising less than 1\% of available GNSS measurements, demonstrating that sub-decimetre localisation is attainable in repetitive row-crop environments without continuous RTK coverage.

    \item A precise-alignment pipeline that registers live point cloud data against a pre-built 3D Gaussian Splatting reference model via Generalised Iterative Closest Point, achieving centimetre-scale pose refinement (15.51\,mm mean inlier RMSE for stereo; 33.25\,mm median inlier RMSE for LiDAR) from a coarse SLAM initialisation.

    \item A comparative evaluation of LiDAR and stereo vision modalities for vine-scale registration, characterising their convergence--accuracy trade-off: LiDAR achieves 100\% condition-level convergence while stereo produces lower RMSE (78\% convergence), with both modalities confirming that the SLAM-derived pose lies within the registration algorithm's convergence basin.
\end{enumerate}